\chapter{Methods of Proof}

In this chapter, we will discuss various methods of proof and provide illustrations of each.  In the first chapter, we saw some examples of proofs without dwelling on the process of doing the proof itself.  The emphasis in this book is to gain experience with reading and writing proofs within various topics in pure mathematics.  It is important to both read \textit{and} write proofs in order to become comfortable with them.  The exit goal of this chapter is that the reader will have a general knowledge of the various types of proofs and be able to read and write such proofs.\\

One may see proofs that are written as lists, say with the statement on the left and the justification on the right, or written more like prose, where more usual written language is combined with the mathematics and logic.  In the previous chapter, there were examples of each type.\\

Much of the material in this chapter is based on the YouTube video series ``Math Major Basics" by MathDoctorBob (https://www.youtube.com/playlist?list=PLF2DF6C3C8015DF5F).\\

We first remark that a \textit{theorem} is a set of valid statements chained together by logical ``ands" (i.e. $\land$).  If all the statements are true and the conclusion is true, then the theorem is \textit{proved}.

\section{Induction}

In the \textit{induction} type of proof, the statements in question will be associated with an index which is used to construct the proof.

\subsection{Weak Induction}

\subsection{Strong Induction}

\section{Proof by Contradiction}

Proof by contradiction is a very useful method of proof which is used quite frequently.  The usual pattern is one assumes the converse of what is being asserted, and then following a sequence of logical deductions arriving at a contradiction.  Once this contradiction has been uncovered, then it implies the original assumption must have been false thus completing the proof.  Here is an example to illustrate the points.


\textit{Prove that there is no rational number whose square is $12$.}  (Note:  This was originally an exercise in Rudin PMA 3rd Edition).

Proof:

\textbf{Suppose that there was a rational number whose square was $12$.}  This means that
$$ \left(\frac{p}{q}\right)^{2}=12$$ for integers $p,q.$  We assume that $p$ and $q$ are not both even.  In other words, all common factors of $2$ have been cancelled.  This can always been done for a rational number.

Upon expansion we obtain:

$$p^{2}=12q^{2}=2\left(6q^{2}\right)$$

This implies that $p^{2}$ and $p$ are both even.  Thus $q$ is odd.  So we write $p=2m$ for integer $m$ and obtain:

$$\left(2m\right)^{2}=4m^{2}=12q^{2}$$

This implies that 

$$m^{2}=3q^{2}$$

Since $q$ is assumed to be odd, $q^{2}$ is odd, and $3q^{2}$ is also odd.  Now, $m^{2}$ must be therefore be odd, and likewise for $m$.  Therefore, for integers $r,n$ we can rewrite the above as:

$$
\left(2r+1\right)^{2}=3\left(2n+1\right)^{2}
$$

Now, expanding both sides yields:
$$
4r^{2}+4r+1=3\left(4n^{2}+4n+1\right)
$$

$$
4r^{2}+4r+1=12n^{2}+12n+3
$$

$$
4r^{2}+4r=12n^{2}+12n+2
$$

$$
4r^{2}+4r=12n^{2}+12n+2
$$

The left side is clearly divisible by 4.

$$
r^{2}+r=3n^{2}+3n+\frac{2}{4}
$$

The left side is an integer, but the right side is not, which is a contradiction. Therefore our assumption about both $p,q$ not even is false.  However, since this can always be done for the rational number $p/q$ it implies that there is no rational $p/q$ such that $\left(p/q\right)^{2}=12.$


\section{Direct Proof}